前言 数学踏上新的征程

人们常说，刚刚过去的一个世纪是数学真正的黄金时代。数学在 20 世 纪的进步比先前所有的世纪加起来还要大。然而，刚刚开始的这个世纪 也可能同样是数学发展的好时候。或许，数学在这个世纪的变迁会和 20 世纪一样巨大，甚至更为惊人。引发这种想法的信号之一是一场渐 变:自 20 世纪 70 年代开始，数学方法的基石---证明的概念逐渐发生 演变，让一个古老却有些被人忽视的数学概念重新回到了舞台中央，这 就是“计算”。

计算能成为引发革命的导火索，这看起来有点不合常理。算法，比如做 加法和做乘法的算法，常常被视为数学知识中最基础的一部分，做计算 也经常被看成是缺乏创造性的枯燥工作。数学家们自己对计算也颇有成 见，勒内·托姆就曾说过:“我的论述中很大一部分属于纯粹的猜想，大 家基本上可以把它们看成是梦话。我接受这种定性......如今，世界上到 处有这么多学者在做计算，难道有人做梦不是件好事吗?”用计算来做 梦，大概还真有点难度啊......

不幸的是，对计算的偏见恰恰根植于“数学证明”这一概念的定义里。确 实，欧几里得以降，“证明”的定义就是利用公理和演绎规则构建的一套 推理 。然而，要解决一个数学问题，仅仅需要构建一套推理吗?数学 的实践难道没有告诉我们，解决问题需要把推理的步骤和计算的步骤巧 妙地融合起来吗?公理化方法若局限在推理中，它所展现的数学视野恐 怕也会十分狭隘。正是因为人们对约束过多的公理化方法多有批评，才 让计算有机会重新出现在数学的舞台上。现在，已有一些研究工作(它 们之间未必有关联)渐渐开始质疑推理高于计算的优势地位，并倡导一 种更为平衡的观点，让两者互为补充。

这场革命让我们重新考量推理和计算之间的关系，同时也促使我们重新 审视数学与物理学、生物学等自然科学之间的对话，特别是数学为何能 在这些学科中发挥难以理解的强大作用这一古老问题，以及自然理论的 逻辑形式这一全新问题。此外，这场革命给“分析判断”和“综合判断”等 哲学概念带来了新的火花。它还让我们反思数学与计算机科学之间的关 系，而且数学似乎是唯一一门不需要借助机器的科学，它为什么如此独 特?

最后，最振奋人心的是，这场革命让我们隐约看到了一些解决数学问题 的新方式，它摆脱了过去的技术强加给证明长度的枷锁---数学也许正 踏上新的征程，去探索从未涉足的全新领域。

诚然，公理化方法的危机并不是凭空出现的。从 20 世纪上半叶起就有 许多先兆，特别是两种理论---可计算性理论和构造性理论的出现。这 两种理论本身虽然没有质疑公理化方法，却重新确立了计算在数学大厦 中的地位。在讨论公理化危机之前，我们会简要回顾这两个概念的历 史。不过，还是让我们先上溯远古，探寻计算这一概念的起源，看看古 希腊人对数学的“发明”过程吧。


##**目标读者**
本书面向学生（本科生或研究生）、工程师和研究人员，他们希望扎实掌握深度学习的实用技术。因为我们从头开始解释每个概念，所以不需要过往的深度学习或机器学习背景。全面解释深度学习的方法需要一些数学和编程，但我们只假设读者了解一些基础知识，包括线性代数、微积分、概率和非常基础的Python编程。此外，在附录中，我们提供了本书所涵盖的大多数数学知识的复习。大多数时候，我们会优先考虑直觉和想法，而不是数学的严谨性。有许多很棒的书可以引导感兴趣的读者走得更远。BelaBollobas的《线性分析》([Bollobás,1999](https：//zh-v2.d2l.ai/chapter_references/zreferences.html#id13))对线性代数和函数分析进行了深入的研究。([Wasserman,2013](https：//zh-v2.d2l.ai/chapter_references/zreferences.html#id178))是一本很好的统计学指南。如果读者以前没有使用过Python语言，那么可以仔细阅读这个[Python教程](http：//learnpython.org/)。
-数学思维与人工智能的关系
-从数学到智能的发展历程